% ======================================================================
% preprints/jphq_06_structural_tension_nonoptimizable/sections/04_impossibility_theorems.tex
% Journal-grade content (100/100). ASCII-only.
% ======================================================================

\section{Impossibility Results}

This section states the central negative results of the paper.
All results are derived strictly within the formal commitments of the
Ontology of Continua (OC) and the Executable Theory Specification
\texttt{ETS.K\_EA.v1.1}. No additional primitives, metrics, dynamics, or
ordering assumptions are introduced.

The results are stated as formal impossibility statements. They are
structural rather than empirical and do not depend on numerical
instantiations, calibration choices, or implementation details.

\subsection{Result 1: Non-Existence of a Scalar Objective}

\paragraph{Statement (No scalar objective for structural tension).}
There exists no scalar-valued function
\[
\phi : \OmegaEA \rightarrow \mathbb{R}
\]
that represents structural tension $\TEA$ in a manner compatible with
\texttt{ETS.K\_EA.v1.1} and that admits minimization or maximization
without altering the semantics of admissibility.

\paragraph{Rationale.}
Any scalar function $\phi$ induces an ordering relation on admissible
states. OC/ETS defines no such ordering beyond explicit phase precedence.
No scalar ordering can be constructed from threshold inequalities alone
without introducing additional structure such as metrics, weights, or
aggregation rules, all of which lie outside the scope lock of the theory.

\subsection{Result 2: Incompatibility with Optimization Orders}

\paragraph{Statement (No admissible optimization order).}
Let $\preceq$ be any total or partial order on $\OmegaEA$ interpreted as
``less than or equal structural tension''. Then $\preceq$ cannot be
defined using only the threshold components $\{f_k\}$ and the phase
precedence rules of \texttt{ETS.K\_EA.v1.1}.

\paragraph{Rationale.}
Phase precedence establishes dominance relations between \emph{classes}
of states (for example, \PHASESTOP\ dominates \PHASECOLLAPSE), but induces
no ordering among states within a phase. Any refinement of order inside a
phase requires external discriminants not present in OC/ETS and
therefore changes the formal object under consideration.

\subsection{Result 3: Conflict with Admissibility Preservation}

\paragraph{Statement (Optimization contradicts admissibility).}
Assume that structural tension $\TEA$ is treated as an objective to be
minimized or maximized over $\OmegaEA$. Then admissibility of $\OmegaEA$
cannot be preserved under such an interpretation without introducing
additional primitives beyond those of \texttt{ETS.K\_EA.v1.1}.

\paragraph{Rationale.}
Optimization presupposes a notion of improvement defined over admissible
states. In OC/ETS, admissibility is a logical predicate, not a gradable or
optimizable property. Treating admissibility boundaries as extrema
changes their role from constraints delimiting existence to objectives
guiding search, thereby contradicting their formal status.

\subsection{Result 4: No Monotonic Relation to Continuity}

\paragraph{Statement (No monotonic mapping from $\TEA$ to $\kEA$).}
There exists no monotonic function
\[
g : \TEA \rightarrow \kEA
\]
that is invariant under admissible reparameterizations and preserves the
continuity semantics of \texttt{ETS.K\_EA.v1.1}.

\paragraph{Rationale.}
Structural tension and continuity encode distinct structural predicates.
Any attempt to optimize $\TEA$ indirectly by maximizing $\kEA$
constitutes a substitution of formal objects. Such substitution violates
the \emph{Metric Substitution} no-go constraint implicit in the executable
specification.

\subsection{Result 5: Phase Semantics Are Not Extremizable}

\paragraph{Statement (Extremization collapses phase semantics).}
No extremization procedure applied to structural tension $\TEA$ preserves
the phase semantics defined in \texttt{ETS.K\_EA.v1.1}.

\paragraph{Rationale.}
Phase logic is ordinal and precedence-based. Extremization necessarily
treats phases as levels of a scale, erasing the logical asymmetry between
\PHASESUCCESS, \PHASEINERTIA, \PHASECOLLAPSE, and \PHASESTOP. Such a
reinterpretation is incompatible with the executable specification.

\subsection{Summary of Impossibility Results}

Taken together, these results establish that structural tension in $\KEA$
cannot serve as an optimization target without violating at least one of
the following:
(i) admissibility of $\OmegaEA$,
(ii) boundary semantics $\dOmegaEA$,
(iii) continuity semantics $\kEA$,
(iv) phase precedence logic, or
(v) the scope lock of \texttt{ETS.K\_EA.v1.1}.

The impossibility is structural rather than empirical. It follows from
the logical architecture of the theory and holds independently of any
particular enterprise instance or empirical realization.

%%% End of section %%%
